\section{Ethics Considerations}
\label{sec:ethics}

This section addresses the ethical dimensions of deploying an
attacker-facing deception system and studying it.

\textbf{Deployment legality and containment.} The honeypot is deployed on
infrastructure the authors own and operate, on a research network reached
through a management plane, with an egress allowlist verified by
configuration audit: the system cannot initiate outbound connections, so
no attacker-controlled content can leave the deception surface. One
predecessor deployment was publicly reachable and interacted with two
genuinely external attackers; those session transcripts were processed
for the aggregate counts reported in \S\ref{sec:advworkload} and not
retained beyond the aggregate. The authors made no institutional-review
filing: the deployment involves no human subjects under the
venue's definition (no human participants were studied, recruited, or
deceived by a human operator), attacker activity on the authors' own
honeypot is the studied phenomenon, and the relevant safeguards are the
containment and privacy measures this section records. The authors will
provide the full configuration-audit trail on request. The
deployment serves AI-generated deception content to parties who connect
to it; the system is nonetheless engineered so that nothing it serves can
contain a real credential: every secret-shape value the fiction can emit
is canary-registered (\S\ref{sec:verify}), making any disclosure a
measured, attributable event rather than a leak. Across the entire
deployment record (279 exported LLM-served sessions, 127 inbound ambient
sessions), zero real credentials left the host and the gate refused 43
natural-traffic drafts fail-closed.

\textbf{Attacker privacy.} All attacker identifiers (IP addresses) are
salted and hashed at the collection boundary (sha256(ip+salt), 12-hex
truncation); raw addresses never leave the collection host, are not
present in any released artifact (verified by redaction audit), and
session transcripts are retained only for the measurement window. No
attacker-identifying information beyond the hashed source address and
shell activity on the honeypot itself is collected or stored.

\textbf{Human subjects.} No human subjects were studied: all attacker
workloads are AI-generated (disclosed throughout), no human raters
participated in any label in this release, and no human-interaction data
beyond the two predecessor-deployment sessions---whose full transcripts
were not retained---is analyzed.

\textbf{Dual-use considerations.} The released artifacts include the
verifier, the fiction-graph schema, and measurement harnesses, but not
the deception content library or canary values; the deception corpus is
required to reproduce the realism numbers and is published only as
redacted previews. The system's defensive purpose (bounding what a
deception endpoint may serve) dominates its offensive surface: the
verifier's misses are published precisely so that defenders can reason
about the residual risk rather than replicate the bypasses.

\textbf{Provider constraints.} The free-tier inference providers used
impose content filters that blocked 41.4\% of replayed attacker commands;
these constraints were recorded and reported, never engineered around
(\S\ref{sec:jeveval}).
